\subsection{skew coneの頂点障害と隣接支持の二固定表示}
\label{subsec:cubic-as4-adjacent-L}
ここでは中心三次が消えない表示を扱う。
因子性や中心の次元だけからadmissibilityを推測せず、切断後のpoint対象の自己Extを調べる。

\begin{lemma}[自己同型型Ore coneの頂点]\label{lem:cubic-skew-cone-vertex}
$B=S[z;\sigma]$ は $\deg z=1$ の自己同型型Ore拡大で、
noetherian quadratic Koszul AS4型14641とする。
$0\ne G\in S_3\cap Z(B)$ が正則であるとし、
$C=S/(G)$, $R=B/(G)=C[z;\bar\sigma]$ と置く。
このとき $P=\pi k[z]\in\qgr R$ は無限高次の自己Extを持ち、$R$ はH3を与えない。
\end{lemma}
\begin{proof}
$S\simeq B/(z)$ はnoetherianであり、Ore基底から
$H_S(u)=(1-u)^{-3}$ となる。
正規正則一次元 $z$ に対するRees lemmaは
$\underline{\operatorname{Ext}}_S^j(k,S)=0$ ($j\ne3$) を与える。
$B$ は左右 $S$ 自由であるから、$M=k\otimes_SB\simeq k[z]$ について
\[
 \underline{\operatorname{Ext}}_B^j(M,B)=0\quad(j\ne3).
\]
中心正則 $G$ のchange of ringsにより
$\underline{\operatorname{Ext}}_R^j(M,R)=0$ ($j\ne2$) となる。
$R$ はAS Gorenstein次元3・parameter 1なので、
定理\ref{thm:h3-polynomial-ore-exclusion}と同じ局所双対とsection比較から
\[
 \operatorname{Ext}^j_{\qgr R}(\pi R(a),P)=0\ (j>0),\qquad
 \operatorname{Hom}_{\qgr R}(\pi R(a),P)=k
\]
が全 $a$ で成り立つ。
$z$ は中心でなくてよい。$M$ 上では $S_{\ge1}$ が零なので、
$z$ 倍写像 $M\to M(1)$ は $R$ 線型であり、有限次元cokernelを持つ。
従って $P\simeq P(1)$ となり、十分高次数のsection次元1から全shiftのHom値を得る。

$C$ 上の $k$ の最小graded自由解消のBetti数を $\beta_n$ とする。
$R$ は左 $C$ 自由なので、これを $-\otimes_CR$ で誘導すると $M$ の自由解消になる。
$\qgr$ に移した各項は上の $P$ に対するHom-acyclicityを持つ。
最小微分の係数は $C_{\ge1}$ に入り $P$ に零作用するから、
first-quadrantのExt spectral sequenceによって
\[
 \dim\operatorname{Ext}^n_{\qgr R}(P,P)=\beta_n
\]
となる。
一方、正則な $G$ とHilbert級数から
\[
 H_C(u)=\frac{1-u^3}{(1-u)^3},\qquad
 \frac1{H_C(u)}=\frac{(1-u)^2}{1+u+u^2}
\]
を得る。後者は多項式でないので、$k$ の最小graded自由解消は有限長ではあり得ない。
従って $\beta_n>0$ は全 $n$ で成立する。
H3ではSerre作用 $\mathcal S=\tau^{-1}[2]$ により $n>2$ のExtが零になるため矛盾する。
これはOre開集合のglobal dimensionを未証明のまま仮定する議論ではない。
\end{proof}

\begin{corollary}[非自明なscalar-cubeの排除]\label{cor:cubic-nontrivial-scalar-cube}
上の自己同型型Ore包絡 $B=S[z;\sigma]$ がCYであり、$S$ は一次生成とする。
$\sigma^3|_{S_1}=\zeta I$, $\zeta\ne1$ なら、$B$ はH3包絡にならない。
\end{corollary}
\begin{proof}
H3を仮定すると、正則中心三次 $G$ が必要である。
Ore基底で $G=\sum_{j=0}^3 f_{3-j}z^j$ と書く。
$z$ との可換性から $\sigma(f_d)=f_d$、従って
$(\zeta^d-1)f_d=0$ である。
$\zeta\ne1$, $\zeta^2\ne1$, $\zeta^3\ne1$ なら $G$ は $kz^3$ に限られるが、
$z^3v=\zeta vz^3$ により非零中心三次にはならない。
$\zeta=-1$ なら $G=f_2z+\alpha z^3$ であり、一次式 $z$ と二次式の積になるため
商は整域でない。
最後に $\zeta$ が非自明三乗根なら $G=f_3+\alpha z^3$。
$S_1$ との可換性を $z$ 次数ごとに比較すると $\alpha=0$ となる。
補題\ref{lem:cubic-skew-cone-vertex}がこの純base切断を排除する。
\end{proof}

\begin{proposition}\label{prop:cubic-adjacent-L}
\cite{AS4lowpinned}の隣接共有支持AZ2にある二表示 $L^\eta_\omega$、
$\omega^2+\omega+1=0$, $\eta=\pm1$ は、どちらもH3包絡にならない。
\end{proposition}
\begin{proof}
具体的な係数は
\[
\begin{gathered}
 \lambda=\omega,\quad x_1=x_4=1,\quad
 x_6=x_{17}=\eta(\omega+2)/6,\quad x_{11}=\eta(\omega+2)/12,\\
 x_8=\eta(\omega+1)/4,\quad x_{15}=\eta/4,\quad
 x_{13}=-1/12,\quad x_{19}=\omega/12
\end{gathered}
\]
で、他の $x_j$ は零である。
生成元 $c$ を除いた $a,b,d$ の三関係で定義される代数を $S$ とする。
残る三関係は厳密に $cv=\sigma(v)c$ の形になり、順序 $(a,b,d)$ で
\[
 \sigma_+=\begin{pmatrix}
 0&0&2(\omega-1)\\-(\omega+2)/6&0&0\\0&\omega&0
 \end{pmatrix},\qquad
 \sigma_-=\begin{pmatrix}
 0&2(1-\omega)&0\\0&0&\omega^2\\(\omega+2)/6&0&0
 \end{pmatrix}
\]
となる。これらが $S$ の関係空間を保つことと可逆性を直接確認したので、
$L^\eta_\omega=S[c;\sigma_\eta]$ は実際のOre表示である。
さらに
\[
 \sigma_+^3=\omega I,\qquad \sigma_-^3=\omega^2I.
\]
元の最高tensor $R A R^{\mathsf T}$ の非零反巡回性も数体上で独立に確認した。
従ってAS4実現ならCYであり、系\ref{cor:cubic-nontrivial-scalar-cube}が両方を排除する。

独立な中心行列計算では、両表示の中心三次は実際に各一次元である。
従ってここは候補中心の消滅による排除ではない。
全候補が $S_3$ に入ることと頂点の無限自己Extを用いており、
三次元の因子性は仮定していない。
\end{proof}

Ore表示・三乗のscalar値・最高tensor・中心三次の一本性と明示生成元は
\texttt{verify\_\allowbreak{}cubic\_\allowbreak{}as4\_\allowbreak{}four\_\allowbreak{}simple\_\allowbreak{}adjacent\_\allowbreak{}L.py}
で検算する。無限自己Extの結論は上の証明に依存し、有限計算そのものの出力と混同しない。
AZ2の具体表示のみを参照し、全chart/CY被覆・AS4存在十分性・全同型は輸入していない。
